\scanpage{97}
\mspage{16}
$x_0$ de $\mathcal C$, mais sans supposer maintenant $\widetilde{\mathcal C}$ connexe. (\deleted{\unsure{En}} Revenant au cas où $\widetilde{\mathcal C}$ proviendrait d'une situation géométrique $X'\mathbin{-}X$, avec $X$ connexe, \deleted{et un \unsure{revêtement}} mais $X'$ seulement $K$-connexe i.e. connexe en tant que schéma sur $K$ --- ce qui signifie qu'il \deleted{\unsure{p}} se décrit en termes galoisiens par $E_0=X'_{x_0}$, sur lequel le produit $\frac12$ direct $\Gamma.\Pi_0$ ($\Pi_0=\Pi_1(X,x_0)=\Pi_1(\mathcal C,x_0)$) opère transitivement, mais sans que l'opération de $\Pi_0$ soit forcément transitive.) Associons à $\forall\ x'\in E_0$ son stabilisateur dans $\Pi_0$, on trouve une application
\diagram{\((50)\qquad
\begin{tikzcd}[ampersand replacement=\&,column sep=large,row sep=small]
E_0 \arrow[r] \& \text{Ens des ss-groupes de }\Pi_0 \\
x' \arrow[r,mapsto] \& \text{stabilisateur }\Pi_{\unsure{0x'|x'}}
\end{tikzcd}\)}
qui commute aux actions de $\Pi_0$ (op. sur le $2^e$ ens par autom. intérieurs), d'où par passage aux quotients
\diagram{\((51)\qquad
\begin{tikzcd}[ampersand replacement=\&,column sep=large]
E_0/\Pi_0\overset{\mathrm{def}}=\{\pi_0\}(E_0) \arrow[r,"{\overline{\Sigma}:}"{pos=.85}] \& \operatorname{Ssgr}(\Pi)/\Pi
\end{tikzcd}\)}
\inserted{$\unsure{\deleted{\pi}|\#}$ c'est le $\pi_0$ dans le topos des $\Pi$-ens.}

D'autre part, $\Theta$ opère sur le $2^{\text{ième}}$ mb de (51) par transport de structure, via l'autom donné de $\Pi$, et \unsure{la|\deleted{la}} une condition nécessai-
\margin{96}
